% TOP_PROBLEM: {"id": 3000058, "title": "Opposite vertices of base polyhedra", "category_id": 2, "category": {"id": 2, "name": "combinatorics", "display_name": "Combinatorics", "description": "Counting problems, graph theory, discrete structures.", "slug": "combinatorics", "order_index": 2, "created_at": "2026-07-31T15:26:25.671Z"}, "status": "open", "problem_number": "AMR-029-0058", "set_id": 13}
% FIRST_POSTED: 2026-10-01
% ENTRY_KIND: statement_only
% UnsolvedMath ID 3000058; source catalogue code AMR-029-0058
% !TeX program = lualatex
% Definition and sources only; this file is not an LLM solution attempt.
\documentclass[11pt,a4paper]{article}
\usepackage[margin=25mm]{geometry}
\usepackage{fontspec}
\setmainfont{lmroman10-regular.otf}[BoldFont=lmroman10-bold.otf,
 ItalicFont=lmroman10-italic.otf,BoldItalicFont=lmroman10-bolditalic.otf]
\usepackage{amsmath,amssymb,mathtools}
\usepackage{unicode-math}
\setmathfont{latinmodern-math.otf}
\AtBeginDocument{\let\setminus\smallsetminus}
\usepackage{xurl,hyperref}
\hypersetup{colorlinks=true,urlcolor=blue}
\setcounter{secnumdepth}{0}
\setlength{\parindent}{0pt}
\setlength{\parskip}{5pt}
\setlength{\emergencystretch}{3em}
\begin{document}
\section*{Opposite vertices of base polyhedra}\label{3000058}
{\small UnsolvedMath ID 3000058 · AMR-029-0058 · Combinatorics · AMR Open Problem Lists}
\subsection{Definitions and mathematical statement}
Let $V=\{1,\ldots,n\}$ be finite and let $f:2^V\to\mathbb R$ be
normalized and submodular:
\[
 f(\varnothing)=0,\qquad
 f(X)+f(Y)\ge f(X\cap Y)+f(X\cup Y)\quad(X,Y\subseteq V).
\]
For $x\in\mathbb R^V$, write $x(X)=\sum_{i\in X}x_i$.
The associated base polyhedron is
\[
            B(f)=\{x:x(X)\le f(X)\ (X\subseteq V),\quad x(V)=f(V)\}.
\]
It is a bounded polyhedron in this finite setting. A vertex is an
extreme point, one that cannot be expressed as a nontrivial convex
combination of two other points of the polyhedron.

Suppose $0\in B(f)$ and every vertex of $B(f)$ belongs to
$\{-1,0,1\}^V$. Must there be a vertex $v$ such that $-v$ is also
a vertex?

The hypothesis $0\in B(f)$ in particular forces $f(V)=0$ and
$f(X)\ge0$ for every $X$. Central symmetry of the entire polyhedron
is not assumed or requested; a single opposite pair is sufficient.
The degenerate polyhedron $\{0\}$ satisfies the conclusion with
$v=0$. Merely expressing the origin as a convex combination of
vertices is weaker: the question asks for an exact pair of additive
opposites among the vertices themselves.
\subsection{Short English statement}
Must a base polyhedron containing the origin and having only
$-1,0,1$ vertices contain a pair of opposite vertices?
\subsection{Sources}
\begin{itemize}
\item[S1] ULAM.AI and the UnsolvedMath Contributors, supplied problems.json, record 3000058 (AMR-029-0058), AMR Open Problem Lists. Catalogue identity and source transcription; CC BY 4.0.
\url{https://huggingface.co/datasets/ulamai/UnsolvedMath}.
\item[S2] Egres Open - List of open problems, Open-problem page: Opposite vertices of base polyhedra. Original source indexed by AMR-029-0058.
\url{https://oldlemon.cs.elte.hu/egres/open/List_of_open_problems}.
\item[S3] EGRES Open, Opposite vertices of base polyhedra. Individual source page and explanatory remarks.
\url{https://lemon.cs.elte.hu/egres/open/Opposite_vertices_of_base_polyhedra}.
\end{itemize}
\end{document}
